\begin{exercisechunk}
\item \textbf{Hard and graded reverse KL.}\suppmark
\label[exercise]{ex:l22-reverse-kl-exactness}
\exerciseprofile{2}{2}{3}{\exproof\quad\excalculation\quad\exinterpretation}{%
Identify why hard reverse KL supplies no finite graded comparison among
imperfect rules and what additional information makes one possible.}
Let $\mathcal U$ be finite, let $\mu(u)>0$ for every $u\in\mathcal U$, and
let $f:\mathcal U\to\{0,1\}$. Define
\[
 P_f(u,y)=\mu(u)\mathbb I\{y=f(u)\},
 \qquad
 Q_q(u,y)=\mu(u)q(y\mid u),
\]
where $q(\,\cdot\mid u)$ is a conditional mass function.
\begin{enumerate}[(a)]
\item Prove that $D(Q_q,P_f)<+\infty$ if and only if
$q(f(u)\mid u)=1$ for every $u\in\mathcal U$. Determine the value of the
divergence when it is finite.
\item Prove
\[
 D(P_f,Q_q)=\sum_{u\in\mathcal U}\mu(u)
 \log\frac{1}{q(f(u)\mid u)}.
\]
Fix $u_0\in\mathcal U$ and $0<\eta<1$. Let
$q(1-f(u_0)\mid u_0)=\eta$ and let $q$ return the correct answer with
probability one at every other prompt. Compute both KL divergences.
\item For fixed $\eta$, explain how the forward divergence can be made
arbitrarily small by changing only $\mu(u_0)$. Explain why the same
construction never satisfies the reverse-KL support requirement.
\item Fix $0<\epsilon<1/2$ and define
\[
 P_{\epsilon,f}(u,y)
 =\mu(u)\bigl((1-\epsilon)\mathbb I\{y=f(u)\}
              +\epsilon\mathbb I\{y=1-f(u)\}\bigr).
\]
For $e_q(u)=q(1-f(u)\mid u)$, prove
\[
 D(Q_q,P_{\epsilon,f})
 =\sum_{u\in\mathcal U}\mu(u)\left(
 (1-e_q(u))\log\frac{1-e_q(u)}{1-\epsilon}
 +e_q(u)\log\frac{e_q(u)}{\epsilon}\right).
\]
Show that the contribution at $u$ diverges as $\epsilon\downarrow0$ when
$e_q(u)>0$. Explain why smoothing makes the objective finite and graded but
does not reveal $f(u)$ at an unobserved prompt.
\item Let $r(\,\cdot\mid u)$ be a positive reference mass function, let
$c:\mathcal U\times\{0,1\}\to\mathbb R$ be an evaluable cost, and let
$\beta>0$. Define
\[
 P_{\beta,c}(u,y)
 =\mu(u)\frac{r(y\mid u)e^{-\beta c(u,y)}}{Z_\beta(u)},
 \qquad
 Z_\beta(u)=\sum_y r(y\mid u)e^{-\beta c(u,y)}.
\]
Prove
\[
 D(Q_q,P_{\beta,c})
 =\sum_{u\in\mathcal U}\mu(u)\left(
 D(q(\,\cdot\mid u),r(\,\cdot\mid u))
 +\beta\sum_yq(y\mid u)c(u,y)+\log Z_\beta(u)\right).
\]
Now take $c(u,y)=\mathbb I\{y\ne f(u)\}$. Show that the minimizing
conditional mass function assigns positive probability to the incorrect
answer at every input for every finite $\beta$. Conclude that the graded Gibbs
objective is not by itself an exact-learning guarantee.

Identify the information, beyond finite positive demonstrations, that permits
an optimization method to evaluate the $q$-dependent terms on generated
prompt--answer pairs.
\end{enumerate}
\begin{hints}
A KL term is infinite when its first argument assigns positive mass to an
event to which its second argument assigns zero mass. For the two graded
targets, substitute their mass functions into the definition of KL and cancel
the common prompt mass $\mu(u)$.
\end{hints}
\end{exercisechunk}
